\documentclass[11pt]{article}
\usepackage[a4paper,margin=2cm]{geometry}
\usepackage{amsmath,amssymb}
\usepackage[hidelinks]{hyperref}
\allowdisplaybreaks
\title{Qpu.Processing.QasmBridge}
\author{Tsvetan Rouschev}
\date{@uuidna/qpu v1.1.1, CC-BY-NC-ND-4.0}
\begin{document}
\maketitle
\noindent\url{https://github.com/uuidna/qpu/blob/v1.1.1/lean/Qpu/Processing/QasmBridge.lean}

\paragraph{\texttt{QasmBridge.registers\_flatten}}
The registers flatten without collision: two qubits laid out at the same position are the same qubit of the same register. The position is `finSigmaFinEquiv`'s, which Mathlib proves to be the sizes of the registers before plus the index (`finSigmaFinEquiv\_apply`). An equivalence is injective (`Equiv.injective`), so equal images mean equal qubits (`Function.Injective.eq\_iff`), and equal positions are equal elements of `Fin` (`Fin.ext\_iff`).
\[ \forall \{m : \mathbb{N}\} \{n : \mathrm{Fin}\ m \to \mathbb{N}\} (a\ b : (i : \mathrm{Fin}\ m) \times \mathrm{Fin} (n\ i)), \uparrow (\mathrm{finSigmaFinEquiv}\ a) = \uparrow (\mathrm{finSigmaFinEquiv}\ b) \to a = b \]
\noindent axioms: \texttt{propext}, \texttt{Classical.choice}, \texttt{Quot.sound}; proof: \url{https://github.com/uuidna/qpu/blob/v1.1.1/lean/Qpu/Processing/QasmBridge.lean\#L17}
\end{document}
